\insertImage{a1837c37-748e-4d37-8201-7339d1d76a59.png}{Aquarelle représentant des leaders aux côtés de leurs employés dans un environnement de travail collaboratif.}

\chapter{Le rôle des leaders dans les entreprises libérées}

Lorsqu'on parle d'entreprises libérées, on a souvent l'impression que la hiérarchie disparaît complètement. Cependant, même si les entreprises libérées cherchent à aplanir la hiérarchie et à décentraliser le pouvoir, cela ne signifie pas que le leadership n'est plus nécessaire. Au contraire, le rôle du leader dans une entreprise libérée est crucial, mais il peut être très différent de ce que nous voyons dans des organisations plus traditionnelles.

Dans ce chapitre, nous allons explorer le rôle des leaders dans les entreprises libérées, comment ils contribuent à la réussite de l'entreprise, les défis auxquels ils sont confrontés et les récompenses qu'ils peuvent récolter. Nous partagerons également des témoignages et des études de cas de leaders qui ont réussi à naviguer dans le monde complexe et parfois déroutant de l'entreprise libérée.

L'objectif de ce chapitre est de montrer que, bien que les entreprises libérées puissent nécessiter un style de leadership différent, elles n'éliminent pas le besoin de leaders forts et efficaces.

\section{Redéfinir le leadership dans une entreprise libérée}

Dans une entreprise libérée, le rôle du leader est profondément transformé. Au lieu de contrôler et de diriger, les leaders sont là pour servir de catalyseurs et de facilitateurs. Ils créent un environnement où les employés peuvent s'épanouir, apporter leurs propres idées et prendre des initiatives.

Dans ce contexte, le leadership ne se limite pas à une position hiérarchique. Il s'agit davantage d'une compétence qui peut être exercée par n'importe qui, à n'importe quel niveau de l'organisation. C'est ce qu'on appelle parfois le "leadership distribué"\footnote{Bolden, R. (2011). Distributed leadership in organizations: A review of theory and research. International Journal of Management Reviews, 13(3), 251-269.}.

Cela signifie que le rôle des leaders formels - comme les managers et les dirigeants - change considérablement. Au lieu de prendre toutes les décisions et de dicter ce que les autres doivent faire, ils sont là pour soutenir et habilité leurs collègues. Ils servent de guides et de mentors, aidant les autres à développer leurs propres capacités de leadership.

Par exemple, chez Michelin, un fabricant de pneus français qui a adopté le modèle d'entreprise libérée, les managers ne sont plus considérés comme des chefs, mais comme des "leaders-coachs". Leur rôle est de favoriser l'autonomie et l'engagement des employés, en fournissant un soutien et des conseils, plutôt que des ordres\footnote{Getz, I. (2009). Liberation management: necessary disorganization for the nanosecond nineties. Michelin.}.

Dans l'ensemble, le leadership dans une entreprise libérée est moins une question de pouvoir et de contrôle, et plus une question de service et d'autonomisation.

\section{Défis et récompenses du leadership libéré}

Devenir un leader dans une entreprise libérée peut être à la fois gratifiant et difficile. Commençons par les défis.

Les leaders dans les entreprises libérées doivent souvent lutter contre des décennies de conditionnement sur ce que signifie être un leader. Ils doivent abandonner l'idée que le leader est celui qui contrôle et décide de tout. Cela peut être déstabilisant et même effrayant pour certains.

De plus, en tant que facilitateur plutôt que décideur, le leader doit souvent faire face à des dilemmes et des paradoxes. Par exemple, comment intervenir dans un conflit sans prendre le contrôle ? Comment aider une équipe à atteindre ses objectifs sans imposer sa vision ? Ces questions ne sont pas faciles à résoudre et nécessitent une grande dose de sagesse et d'humilité.

Malgré ces défis, être un leader dans une entreprise libérée peut aussi être extrêmement gratifiant.

D'abord, il y a la satisfaction de voir les employés s'épanouir et prendre des initiatives. De nombreux leaders rapportent qu'il est très gratifiant de voir leurs collègues grandir et s'épanouir dans un environnement de travail libéré\footnote{Laloux, F. (2014). Reinventing organizations: A guide to creating organizations inspired by the next stage of human consciousness. Nelson Parker.}.

Ensuite, il y a le plaisir de contribuer à quelque chose de plus grand que soi. Les entreprises libérées cherchent souvent à avoir un impact positif non seulement sur leurs employés et leurs clients, mais aussi sur la société dans son ensemble. En tant que leader, vous pouvez jouer un rôle clé dans la réalisation de cette vision.

Enfin, il y a la possibilité d'apprendre et de grandir en tant que personne. Le voyage vers le leadership libéré est souvent un voyage de développement personnel, où vous êtes amené à questionner vos croyances et vos habitudes, à apprendre de nouvelles compétences et à devenir une meilleure version de vous-même.

\section{Témoignages de leaders et études de cas}

Pour illustrer le rôle des leaders dans les entreprises libérées, voici quelques témoignages et études de cas de leaders qui ont réussi à naviguer dans le monde complexe et parfois déroutant de l'entreprise libérée.

\subsection{Témoignage 1 : Un leader chez Poult}

Poult, le fabricant de biscuits français que nous avons mentionné précédemment, est dirigé par une équipe de leaders qui se considèrent davantage comme des "facilitateurs" que comme des "chefs". Un de ces leaders partage son expérience :

"En tant que leader chez Poult, mon rôle est de soutenir mes collègues et de leur donner les moyens de réussir. Cela peut être difficile parfois, car je dois résister à l'envie de prendre le contrôle et de résoudre les problèmes moi-même. Mais quand je vois comment les gens s'épanouissent et prennent des initiatives, je sais que cela en vaut la peine"\footnote{Rodriguez, B. (2014). Poult. L'Usine nouvelle.}.

\subsection{Témoignage 2 : Un leader chez FAVI}

FAVI, la fonderie française, est dirigée par une équipe de leaders qui croient fermement à l'autonomie et à la responsabilisation. Un de ces leaders décrit son expérience comme suit :

"Chez FAVI, je suis moins un manager qu'un coach. Mon rôle est d'aider mes collègues à développer leurs compétences et à prendre des décisions. C'est un défi constant, car chaque personne est unique et a ses propres besoins. Mais c'est aussi extrêmement gratifiant de voir les gens grandir et devenir plus confiants et autonomes"\footnote{Getz, I. (2009). Liberation management: necessary disorganization for the nanosecond nineties. Favi S.A.}.

Ces témoignages montrent que le leadership dans une entreprise libérée peut être à la fois un défi et une source d'épanouissement personnel et professionnel.
